\section{Repository Deliverables}
\begin{longtable}{p{0.22\textwidth}p{0.34\textwidth}p{0.35\textwidth}}
\caption{Evidence and implementation boundaries}\label{tab:submission-evidence}\\
\toprule
\textbf{Deliverable} & \textbf{Repository evidence} & \textbf{Boundary}\\
\midrule\endfirsthead
\toprule\textbf{Deliverable} & \textbf{Repository evidence} & \textbf{Boundary}\\\midrule\endhead
Mobile app & \path{apps/mobile/lib/}, \path{apps/mobile/test/} & Flutter source and demonstration data; device build and shared API integration require separate evidence.\\
Web portals & \path{apps/web/src/}, \path{apps/web/package.json} & React/TypeScript interface for lab and admin workflows; local demonstration state.\\
Business API & \path{backend/server.js}, \path{backend/authMiddleware.js} & Express endpoints with PostgreSQL dependency; configured database and integration tests are required.\\
AI prototype & \path{ai_service/}, \path{ai_service/openapi.json} & Local recommendation, retrieval, grounded extraction, and scan-quality heuristics.\\
AI validation & \path{tests/}, \path{scripts/evaluate_ai.py}, \path{scripts/smoke_test.py} & Automated tests and reproducible evaluation scripts; synthetic/local results do not demonstrate production quality.\\
Deployment design & \path{Dockerfile}, \path{infra/azure/}, \path{.github/workflows/} & Container, infrastructure, and workflow definitions; presence of files is not proof of a live deployment.\\
Design report & \path{CNPM_Report.tex}, \path{sections/}, \path{docs/} & Consolidated requirements, responsibilities, architecture, interface flows, and submission evidence.\\
\bottomrule
\end{longtable}

\section{Demonstration Procedure}
\begin{enumerate}
\item Install the web dependencies in \path{apps/web/} with
\texttt{npm ci}; run \texttt{npm run build} and \texttt{npm run dev}.
Open the local address printed by Vite.
\item Demonstrate Film Lab and Admin portal navigation, order search/filtering,
processing stages, marketplace/archive views, and available form validation.
Describe each operation as a local prototype action unless a verified API
response and persistent database result are shown.
\item For mobile, install a compatible Flutter SDK, run
\texttt{flutter pub get}, \texttt{flutter analyze}, and \texttt{flutter test}
inside \path{apps/mobile/}. Launch a supported device target with
\texttt{flutter run}. Source inspection alone is not a device test.
\item From the repository root, install Python project dependencies and run
\texttt{python -m unittest discover -s tests -v}. Start the AI adapter with
\texttt{python -m ai\_service.api}; use \path{/health}, \path{/ready}, and
\path{/openapi.json} to inspect its contract.
\item Present one supported photography question, one unsupported question,
and one recommendation request. Explain the local knowledge/data boundary and
show how insufficient evidence is reported.
\item Open this report and the acceptance/submission documents in \path{docs/}.
Relate each demonstrated feature to the relevant requirement and work item.
\end{enumerate}

\section{Acceptance and Remaining Integration Work}
For each requirement, record its identifier, owner, code or document path,
test procedure, observed result, and unresolved dependency. A feature is
complete only when its agreed acceptance criteria have supporting evidence.
Interface mockups, implemented local interactions, automated unit tests, and
verified end-to-end operations are distinct evidence levels.

Before claiming the complete platform is operational, the team must verify:
\begin{itemize}
\item database schema and repeatable migrations, seeded non-sensitive test
data, configured secrets, and authorization for each protected resource;
\item mobile/web authentication and shared API integration, including
cross-user isolation, failure states, and persistence after reload;
\item private cloud uploads, ownership checks, scan publication and downloads;
\item real payment, logistics, and notification adapters, or an explicitly
documented decision to exclude them from the assessed prototype;
\item supported-device mobile builds, end-to-end tests, accessibility checks,
and performance/recovery measurements against the proposed targets;
\item deployment health, backup and restore verification, and reproducible
release configuration before calling the system production ready.
\end{itemize}

\section{Contribution and Submission Evidence}
The source repository is the primary record of implemented changes. Each
meaningful commit should describe a coherent change and, where available,
reference its Jira issue. Jira items should link to the resulting code,
document, test evidence, or demonstration and describe unresolved dependencies.
Only move an item to Done when its actual acceptance criteria are satisfied.

Contributor graph height is a count-based view of eligible commits in its
selected period. It does not measure feature completeness or the quality of a
member's work. The submission should explain ownership with requirements,
implemented artifacts, tests, and review evidence, without representing planned
work as completed work.
